On considère une solution aqueuse d'acide propanoïque $\ce{C2H5CO2H}$ de volume
$V$, de concentration molaire $C=\qty{2e-3}{\mol\per\L}$ et de
$\mathrm{pH}=3,79$ à $\qty{25}{\degreeCelsius}$.

\begin{enumerate}
    \item Écrire l'équation chimique modélisant la réaction de l'acide
          propanoïque avec l'eau.
    \item Calculer la valeur du taux d'avancement $\tau$ de la réaction.
          Conclure.
    \item Montrer que l'expression de la constante d'acidité $\mathrm{K_{A1}}$
          du couple $(\ce{C2H5CO2H_{(aq)}}/\ce{C2H5CO^-_{2(aq)}})$ s'écrit~:
          $\mathrm{K_{A1}}=\num{1.43e-5}$. Vérifier que
          $\mathrm{K_{A1}}=\frac{10^{-2\mathrm{pH}}}{C-10^{-\mathrm{pH}}}$
    \item Représenter le diagramme de prédominance des deux espèces du couple
          $(\ce{C2H5CO2H_{(aq)}}/\ce{C2H5CO^-_{2(aq)}})$ présentes dans
          la solution étudiée.
    \item On considère l'acide benzoïque de formule $\ce{C6H5CO2H}$. On
          note $\mathrm{K_{A2}}$ la constante d'acidité du couple
          $(\ce{C2H5CO2H_{(aq)}}/\ce{C2H5CO^-_{2(aq)}})$. Pour déterminer la
          valeur de $\mathrm{K_{A2}}$, on mélange le même volume de la solution
          aqueuse d'acide propanoïque et d'une solution aqueuse de benzoate de
          sodium $\ce{C6H5CO^-_{2(aq)} + Na^+_{(aq)}}$. Les deux solutions ont
          même concentration molaire.
          \begin{enumerate}
              \item Écrire l'équation chimique de la réaction qui se produit
                    entre l'acide propanoïque $\ce{C2H5CO2H_{(aq)}}$ et l'ion
                    benzoate $\ce{C6H5CO^-_{2(aq)}}$.
              \item Recopier, sur votre copie, le numéro de la question et
                    écrire la lettre correspondante à la proposition vraie.

                    L'expression de la constante d'équilibre $\mathrm{K}$
                    associée à l'équation chimique de cette réaction est~:
                    \begin{minipage}{\linewidth}
                        \begin{table}[H]
                            \centering
                            \setlength{\arrayrulewidth}{0.8pt}
                            \renewcommand{\arraystretch}{2}
                            \begin{tabularx}{\textwidth}{|c|C|c|C|c|C|c|C|}
                                \hline 
                                \raisebox{1mm}{$\mathbf{A}$} &
                                \raisebox{1mm}{$\mathrm{K}=
                                \frac{\mathrm{K_{A2}}}{\mathrm{K_{A1}}}$} &
                                \raisebox{1mm}{$\mathbf{B}$} &
                                \raisebox{1mm}{$\mathrm{K}=
                                \mathrm{K_{A1}}\cdot\mathrm{K_{A2}}$} &
                                \raisebox{1mm}{$\mathrm{C}$} &
                                \raisebox{1mm}{$\mathrm{K}=
                                \frac{\mathrm{K_{A1}}}{\mathrm{K_{A2}}}$} &
                                \raisebox{1mm}{$\mathrm{D}$} &
                                \raisebox{1mm}{$\mathrm{K}=
                                \frac{1}{\mathrm{K_{A1}}\cdot\mathrm{K_{A2}}}$}\\
                                \hline
                            \end{tabularx}
                        \end{table}
                    \end{minipage}
          \end{enumerate}

    \item Calculer la valeur de $\mathrm{K_{A2}}$ sachant que $\mathrm{K}=0,23$.
\end{enumerate}


